%% ============================================================
%% 数据包络分析 (DEA) 段落模板
%% 变量:
%%   n_dmu                — 决策单元数量
%%   dmu_names            — DMU 名称列表
%%   input_names          — 投入指标名称列表
%%   output_names         — 产出指标名称列表
%%   model_type           — 模型类型 ("CCR" / "BCC")
%%   orientation          — 导向 ("input" / "output")
%%   efficiency_scores    — 效率值列表
%%   is_efficient         — 是否 DEA 有效 (bool 列表)
%%   slacks_inputs        — 投入松弛变量 (二维列表, 可选)
%%   slacks_outputs       — 产出松弛变量 (二维列表, 可选)
%%   efficiency_bar_path  — 效率值柱状图路径
%% ============================================================

\subsection{基于数据包络分析（DEA）的效率评价}

\subsubsection{模型原理}

数据包络分析（Data Envelopment Analysis, DEA）是一种基于线性规划的非参数效率评价方法，无需假设生产函数的具体形式。

本文采用 \textbf{<< model_type >>} 模型（<< "投入导向" if orientation == "input" else "产出导向" >>），共有 $<< n_dmu >>$ 个决策单元（DMU），投入指标 $<< input_names | length >>$ 个，产出指标 $<< output_names | length >>$ 个。

<% if model_type == "CCR" %>
CCR 模型假设规模报酬不变（CRS），对于第 $k$ 个 DMU，效率评价的线性规划对偶形式为：
\begin{equation}
\begin{aligned}
    \min \quad & \theta \\
    \text{s.t.} \quad & \sum_{j=1}^{n} \lambda_j x_{ij} \leq \theta x_{ik}, \quad i = 1, \ldots, s_1 \\
    & \sum_{j=1}^{n} \lambda_j y_{rj} \geq y_{rk}, \quad r = 1, \ldots, s_2 \\
    & \lambda_j \geq 0, \quad j = 1, \ldots, n
\end{aligned}
\label{eq:dea_ccr}
\end{equation}
<% else %>
BCC 模型放松了规模报酬假设（VRS），增加约束 $\sum \lambda_j = 1$，可区分纯技术效率与规模效率。
<% endif %>

当 $\theta^* = 1$ 且所有松弛变量为零时，该 DMU 为 DEA 有效。

\subsubsection{计算结果}

\begin{table}[H]
    \centering
    \caption{<< model_type >> DEA 效率评价结果}
    \label{tab:dea_result}
    \begin{tabular}{lcc}
        \toprule
        \textbf{DMU} & \textbf{效率值 $\theta^*$} & \textbf{是否有效} \\
        \midrule
<% for i in range(n_dmu) %>
        << dmu_names[i] >> & << "%.6f" | format(efficiency_scores[i]) >> & << "是" if is_efficient[i] else "否" >> \\
<% endfor %>
        \bottomrule
    \end{tabular}
\end{table}

<% if efficiency_bar_path %>
\begin{figure}[H]
    \centering
    \includegraphics[width=0.8\textwidth]{<< efficiency_bar_path >>}
    \caption{各 DMU 效率值对比}
    \label{fig:dea_bar}
\end{figure}
<% endif %>

共有 << is_efficient | select | list | length >> 个 DMU 达到 DEA 有效，占比 << "%.1f" | format(is_efficient | select | list | length / n_dmu * 100) >>\%。效率值最低的为 << dmu_names[efficiency_scores.index(min(efficiency_scores))] >>（$\theta^* = << "%.4f" | format(min(efficiency_scores)) >>$），可进一步分析其投入冗余与产出不足。